\section{计算协议}\label{app:computational-protocol}

\subsection{参考环境}

公开证明材料是不可变 Git 快照

\begin{center}
\url{https://github.com/whzy3185/math/tree/bb3c8aca39a27c76e8b76bbe733b84d0fd1fafc6}
\end{center}

相应的投稿材料清单为
\path{research/reproducibility/}\allowbreak\path{target_a_submission_}\allowbreak\path{artifact_manifest.json}。
该清单记录定理~\ref{thm:smallest-counterexample} 和定理~
\ref{thm:low-period-frontier} 的材料角色与 SHA-256 摘要。验证器从固定 Git
对象读取指定文件，而不会默认接受可变工作树副本。该快照包含全部检查器源码、
完整 2,626 行低周期证书表、24 至 30 阶检查点数据、紧凑最小性材料、重新计算
摘要和固定依赖。投稿时最好提供带 DOI 的存档版本；完整提交标识符已经使本稿
的公开材料不可变且按内容寻址。

参考解释器和软件包为：

\begin{center}
\begin{tabular}{ll}
\toprule
组件 & 版本 \\
\midrule
Python & 3.12.13 \\
NumPy & 2.3.5 \\
SymPy & 1.14.0 \\
pytest & 9.1.1 \\
\bottomrule
\end{tabular}
\end{center}

在具有 Python 3.12 的干净检出中，用下列命令创建可移植环境：

\begin{lstlisting}[language=bash]
python3.12 -m venv .venv-target-a
.venv-target-a/bin/python -m pip install --upgrade pip
.venv-target-a/bin/python -m pip install \
  -r research/reproducibility/requirements-target-a.txt
\end{lstlisting}

下文所有命令均相对于仓库根目录。

\subsection{快速验证}

完整默认回归套件运行命令为

\begin{lstlisting}[language=bash]
.venv-target-a/bin/python -m pytest -q research/scripts
\end{lstlisting}

也可分别运行主要定理检查器：

\begin{lstlisting}[language=bash]
.venv-target-a/bin/python research/scripts/verify_target_a_minimality_certificate.py
.venv-target-a/bin/python research/scripts/verify_target_a_period8_infinite_family.py
.venv-target-a/bin/python research/scripts/verify_target_a_period8_sharp_constant.py
.venv-target-a/bin/python research/scripts/verify_target_a_period8_structural_mechanism.py
.venv-target-a/bin/python research/scripts/verify_target_a_general_period_moments.py
.venv-target-a/bin/python research/scripts/verify_target_a_low_period_spectral_frontier.py
.venv-target-a/bin/python research/scripts/verify_target_a_low_period_structural_frontier.py
.venv-target-a/bin/python research/scripts/verify_target_a_periodic_operator_equivalences.py
.venv-target-a/bin/python research/scripts/verify_target_a_submission_artifact_manifest.py
\end{lstlisting}

只要依赖、计数、精确证书、范围标志或源码摘要不匹配，相应命令即终止并报告失败。

\subsection{慢速生成器核验}

最强的代表元集合核验由下列命令直接执行：

\begin{lstlisting}[language=bash]
.venv-target-a/bin/python \
  research/scripts/target_a_record_set_audit.py \
  --n 24 26 28 30
\end{lstlisting}

该程序构建独立 C 扫描器，将正式 FKM 代表元写入磁盘映射表，并要求逐项消耗
整个集合。输出中的缺失、重复、剩余、规范性及轨道大小错误均为零，同时给出
精确缺陷与轨道大小直方图。此项核验重构完整代表元集合，但不重新计算每个状态
的谱判定。

独立的大阶谱判定由下列命令重新执行：

\begin{lstlisting}[language=bash]
.venv-target-a/bin/python \
  research/scripts/target_a_independent_spectral_audit.py \
  --n 24 26 28 30
\end{lstlisting}

该命令读取 C 全空间扫描器导出的记录，独立重构两种和乐矩阵，并用精确整数
Rayleigh 商验证每个非最优状态。它重新执行全部判定，而非只复核检查点。记录
的参考运行耗时为数十分钟；每阶计数、独立证书摘要、源码摘要及零未认证状态
由 \path{research/reproducibility/}\allowbreak
\path{target_a_computational_evidence_manifest.json} 绑定。

默认套件跳过的三个生成器核验需显式启用：

\begin{lstlisting}[language=bash]
TARGET_A_RUN_SLOW_GENERATOR_TESTS=1 \
.venv-target-a/bin/python -m pytest -q \
  research/scripts/test_target_a_direct_bracelets.py::DirectBraceletTests::test_direct_generator_burnside_n26 \
  research/scripts/test_target_a_direct_bracelets.py::DirectBraceletTests::test_direct_generator_burnside_n28 \
  research/scripts/test_target_a_direct_bracelets.py::DirectBraceletTests::test_direct_generator_burnside_n30
\end{lstlisting}

这些测试独立检查壳层总数、Burnside 计数、所代表类总数、次序和奇偶性。
这些测试用于独立核验，而不是重新进行谱计算。

\subsection{重新计算与完整性复核}

新的计算从规范生成器开始，重构每个谱状态，重新计算其精确判定，并写入新的
检查点序列。完整性复核读取已有序列并验证其摘要、计数、游标和汇总记录。二者都
有用，但只有前者会对每个状态重新执行数学判定。

对各正式计算阶数，选择一个空的外部目录并运行

\begin{lstlisting}[language=bash]
export OUT="${TARGET_A_OUTPUT:?set TARGET_A_OUTPUT}"
mkdir -p "$OUT/checkpoints" "$OUT/results"
for n in 24 26 28 30; do
  .venv-target-a/bin/python research/scripts/target_a_minimality_search.py \
    --n "$n" \
    --checkpoint-dir "$OUT/checkpoints/n$n" \
    --output "$OUT/results/target_a_search_n$n.json"
done
\end{lstlisting}

在某阶数中断后，可向其命令添加 \(--resume\) 继续。每个结果 JSON 记录完成
状态和壳层总数、有序输入与证书摘要、末端检查点链摘要、耗时和峰值驻留内存。
在 Apple M5 上使用 Python 3.12.13、NumPy 2.3.5 和 SymPy 1.14.0 的一次
成功参考运行得到：

\begin{table}[tbp]
\centering
\caption{重新计算所用资源与末端检查点摘要。}
\label{tab:regeneration-resources}
\begingroup
\renewcommand{\arraystretch}{1.12}
\scriptsize
\resizebox{\textwidth}{!}{%
\begin{tabular}{lllllll}
\toprule
\(n\) & 状态数 & 分块数 & 时间（秒） & 峰值 RSS（MiB） & 检查点磁盘 & 末端链 SHA-256 \\
\midrule
24 & 353,812 & 28 & 25.3 & 75.1 & 1.1 MiB & \texttt{\detokenize{2dde869aea5da4f040e67a4fef3e93b5f35f5fb42d75f6d48820439f418c83c1}} \\
26 & 1,299,064 & 76 & 117.2 & 84.4 & 3.1 MiB & \texttt{\detokenize{c515350b8bea840c04448086fbc98523615364c05ca837553c93efb933bc0c4e}} \\
28 & 4,810,472 & 250 & 414.9 & 86.7 & 12 MiB & \texttt{\detokenize{7ba200b05590b2a9c0ea1121f25f89ddf1d294d0c043417e69f78f764f6e8ee1}} \\
30 & 17,929,600 & 908 & 1542.3 & 135.2 & 43 MiB & \texttt{\detokenize{b7fd264eece645eead187424152ae810a9ff940e37ffc5649b5ddf65aa31d59d}} \\
\bottomrule
\end{tabular}
}
\endgroup
\end{table}

运行时间依赖硬件；观测到的串行总时间约 35 分钟，单进程峰值内存低于
150 MiB，四个检查点目录共占约 60 MiB。已提交的紧凑再现摘要还记录有序
输入与有序证书摘要，所以无需仓库外计时日志即可比较新的计算结果。仓库中保存
的检查点序列由下列命令复核：

\begin{lstlisting}[language=bash]
.venv-target-a/bin/python research/scripts/target_a_checkpoint_replay.py \
  --n 24 26 28 30 \
  --output "$OUT/target_a_checkpoint_replay.json"
\end{lstlisting}

正式计算的检查点已保存在仓库中。新的计算分块和完整运行日志保存在仓库外；其紧凑清单
记录状态数、分块数、有序输入与证书摘要以及末端检查点序列摘要。复核要求
逻辑摘要一致，而不要求计时元数据逐字节一致；未发现定义不一致。

\subsection{错误注入测试}

错误注入测试每次修改一个证明数据字段，并验证检查程序能够识别错误。所检验的异常
包括遗漏轨道代表
元、重复规范字、非确定性轨道标识符、依赖摘要漂移、错误本原周期、被改动的
Rayleigh 分子、较低根式分支、错误矩方向、把数值预览提升为证明，以及禁止的
全周期范围。此类异常检验不可省略，因为多项正向检查共享同一组精确算术原语。
